\section{Introduction}

Firms that lay off the most workers tend to cut wages the least. Why do some firms lay off workers rather than cut wages? In most of the literature, the answer is that firms cannot cut wages for exogenous reasons and therefore resort to layoffs. The assumption of exogenous wage rigidity—potentially justified by minimum wages, morale costs of cuts, or sectoral bargaining—is pervasive in macroeconomics. It is used to amplify employment fluctuations and firms’ responses to monetary shocks. However, recent survey evidence (\textcite{davis2025}, \textcite{bertheau2025}) suggests that firms often have discretion and sometimes actively prefer layoffs to feasible wage cuts.

In this paper, I offer an alternative explanation for why firms lay off workers instead of cutting wages. Layoffs and wages are instruments of managing their workforce. During negative shocks, firms have two options: fire poor matches or cut wages incentivizing workers to leave. Following recent survey paper by \textcite{bertheau2025}, firms find layoffs to be a more valuable tool: “Layoffs give better control over who leaves the firm.” is the modal answer firms gave to the question of “Why didn’t you lower pay instead of laying off employees?”. I build a dynamic frictional labor market model with this trade-off in mind. Each period, firms hire workers; each match has a quality realized after work begins and observed only by the firm. To facilitate that layoffs are the main instrument for shedding poor matches, the firm cannot wage-discriminate within a hiring cohort (workers hired in the same period). Layoffs are thus the only tool to affect a cohort’s composition. When used, they raise the cohort’s average quality and reduce the firm’s incentive to cut survivors’ wages.


%In this paper, I offer an alternative explanation for why firms lay off workers instead of cutting wages. Rather than lay off because they must, firms \emph{choose} to lay off and then avoid cutting survivors’ wages. For intuition, consider a survey question in \textcite{bertheau2025}: “Why didn’t you lower pay instead of laying off employees?” The modal answer is: “Layoffs give better control over who leaves the firm.” I interpret this as firms finding some workers more productive than others and preferring layoffs as a tool to remove lower-productivity matches (whereas wage cuts could induce valuable workers to exit). I build a dynamic model with a frictional labor market in which both unemployed and employed workers search. Each period, firms hire workers; each match has a quality realized after work begins and observed only by the firm. To facilitate that layoffs are the instrument for culling poor matches, the firm cannot wage-discriminate within a hiring cohort (workers hired in the same period). Layoffs are thus the only tool to affect a cohort’s composition. When used, they raise the cohort’s average quality and reduce the firm’s incentive to cut survivors’ wages.

Beyond surveys, I use French matched employer–employee data from 2009–2019 to discipline and validate the model. First, I confirm the literature’s finding (\textcite{enrlich2024}) that firms which lay off the most exhibit the most rigid wages. I document the layoff rate and the pass-through of firm-level productivity shocks to the wages of job stayers. Grouping firms by average layoff rates, I find that after a positive productivity shock normalized to $100\%$, wages in firms firing the least increase by $1.7\%$, whereas wages in firms firing the most change by less than $1\%$.

Furthermore, the connection between layoffs and wage rigidity holds not only across firms but also within firms, across tenure. Workers with tenure below two years face layoff rates of $5$--$11\%$, compared with $2\%$ for workers with four to five years of tenure. By contrast, wages of workers with less than one year of tenure essentially do not respond to negative productivity shocks, whereas more senior workers’ wages fall by $2.4\%$ in response to a $100\%$ productivity shock.

Lastly, I examine relative wage changes between junior and senior workers when a firm begins layoffs. Whenever firms lay off workers, a wage gap opens between senior and junior contracts.

To understand these facts, the paper builds an equilibrium model of the labor market with search frictions and one-sided limited commitment on the worker side. Firms employ a continuum of workers, exhibit decreasing returns to scale, and face firm-level productivity shocks. Search is directed: firms post dynamic contracts to attract workers. Contracts specify future productivity-history-contingent wages and layoff risk. Workers are risk-averse and cannot commit to stay; firms can credibly insure workers against future productivity shocks. Firms thus face a trade-off between insuring workers against future losses and incentivizing workers to stay (leave) when they are most (least) productive. While ex-ante homogeneous, all matches receive an initial permanent shock at the start of employment. Firm-level productivity is common knowledge; the match-specific shock is observed only by the firm.

Because firms have decreasing returns to scale, they manage contracts jointly across employees. In principle, this renders the model intractable: in a recursive formulation, the firm would need to track contracts for a continuum of workers, producing an infinite-dimensional state space. Fortunately, with directed search, within a given period a firm hires all workers on the same contract. It is therefore sufficient to track contracts by cohort. This discretizes the state space and makes it finite for firms of finite age: a firm of age ten employs at most ten cohorts. Moreover, up to an approximation, the state space can be bounded: I show theoretically and empirically that cohort wages tend to converge. I use this to justify pooling workers beyond a tenure threshold onto the same contract. Thus, instead of optimizing a continuum of contracts, it suffices to work with a finite set of tenure-specific contracts.

I assume firms cannot wage-discriminate by match quality within a cohort. Layoffs are therefore the only way to affect within-cohort quality. In low-productivity states, layoffs serve two purposes: they shed labor that is overpaid relative to productivity and cleanse the workforce of poor matches. In cohorts that experience layoffs, surviving workers are of higher average quality. The firm is then more inclined to retain survivors and cuts their wages less than in cohorts whose composition did not change. This mechanism is central to the paper’s view of layoffs and pay cuts: wherever the firm chooses to fire workers, it has less desire to cut survivors’ wages.

Unlike models with exogenous wage rigidity, this framework explains the layoff–wage connection both between and within firms. When choosing whom to lay off, it is optimal to eliminate poor matches—typically the least costly workers to dismiss, often juniors. The model thus generates relatively higher layoff rates among the most junior workers. Individual layoff risk depends on both absolute tenure and relative tenure within the firm: the most junior (and least costly) workers are cut first, consistent with last-in, first-out patterns documented by \textcite{buhai2014}. After layoffs, the firm has less incentive to cut the wages of surviving juniors than of seniors, whose average quality has not improved. The model can therefore replicate the tenure-gradient in layoffs and wage rigidity documented in the data.

The no–within-cohort discrimination assumption is central. The quoted intuition—“Layoffs give better control over who leaves the firm”—requires that layoffs be more useful than wage cuts for removing undesirable matches. Under complete information about match quality, layoffs would play no special role and the model would not generate meaningful layoffs. I offer several justifications for the assumption. One interpretation is nondiscrimination law: if firms cannot justify why they prefer one worker to another, they cannot selectively change wages. This is distinct from observable differences in worker quality. There is no doubt about productivity differences when Ann publishes significantly more than Bob, and the model need not restrict contracts based on observable quality. It is enough to assume the existence of some firm preferences over workers that are not known to workers themselves.

In the asymmetric-information case, I also consider an extension in which the firm may wage-discriminate. There remains a reason not to: as long as workers do not know their own quality, they effectively “risk-share’’ the layoff risk that arises. This makes layoffs cheaper for the firm, yielding two opposing instruments for removing poor matches. In Appendix~\ref{microfoundation}, I analyze a signaling game in which wages affect workers’ beliefs about their type. Under a sufficiently inelastic probability of workers being retained (given outside opportunities), a pooling equilibrium exists in which the firm does not reveal match types. Focusing on this equilibrium justifies the maintained assumption.

\textbf{Related Literature}
This paper builds on the literature on the roots of wage rigidity. Broadly, two views exist. One treats rigid wages as the endogenous outcome of bilaterally efficient bargaining (\textcite{barro1977}, \textcite{thomas1988}, \textcite{balke2022}). The other models rigid wages as the result of exogenously imposed frictions (\textcite{christiano2005}, \textcite{nekarda2020}, \textcite{blanco2025}). Only the latter can directly connect wage rigidity to layoffs: if firms are restricted from cutting wages, they resort to layoffs. This paper connects wage rigidity and layoffs while taking inspiration from the bilaterally efficient view: firms optimally choose both to fire some workers and to refrain from cutting survivors’ wages.

On the empirical side, recent survey evidence builds on \textcite{bewley1999}. \textcite{bertheau2025} and \textcite{davis2025} are closest to this paper, both studying the wage-cut–layoff trade-off. \textcite{davis2025} survey unemployed workers and find that, while many were open to substantial pay cuts, firms rarely initiated such discussions before resorting to layoffs. Some respondents also conjecture that wage cuts could induce the best workers to leave, suggesting firms do not have complete control over individual wages. On the firm side, \textcite{bertheau2025} merge a large-scale firm survey with administrative data to study incentives to cut wages versus fire workers. They find that wage cuts are often a poor substitute for layoffs because firms want to remove particular workers. This paper provides a theoretical foundation for that intuition.

The model sits at the intersection of dynamic contracts, firm dynamics, and match heterogeneity. 

First, it relates to optimal wage contracts in long-term employment. Building on \textcite{thomas1988} and \textcite{holmstrom1982}, optimal contracts have been studied in rich search environments by \textcite{burdett2003}, \textcite{menzio2011}, \textcite{fukui}, \textcite{balke2022}, \textcite{souchier2022}, \textcite{cedomir2024}, and \textcite{elsby2024}, with emphasis on on-the-job search. Closest are \textcite{balke2022}, \textcite{souchier2022}, and \textcite{cedomir2024}, who study the insurance–incentives trade-off between risk-averse workers and risk-neutral firms. My contribution is to incorporate layoffs into the insurance contract: firms can pass through negative shocks by either cutting wages or firing workers. Within dynamic insurance contracts, this is also the first paper to allow decreasing returns to scale, so firms jointly manage insurance across their workforce. The closest counterpart is \textcite{schaal2017}, where firms with DRS production offer dynamic contracts to risk-neutral workers.

Second, it connects to search-and-matching models with firm dynamics, including \textcite{acemoglu2014}, \textcite{elsby2013}, \textcite{kaas2015}, \textcite{schaal2017}, and more recently \textcite{gulyas2020}, \textcite{bilal2022}, \textcite{elsby2022}, \textcite{mccrary2022}, \textcite{rudanko2023}. Beyond \textcite{schaal2017}, this paper introduces dynamic contracts with worker–firm risk insurance into this setting. Dynamic contracts are crucial to reconciling within-firm heterogeneity in layoffs and wage pass-through: unlike models with Nash bargaining (\textcite{mccrary2022}) or sequential bargaining (\textcite{bilal2022}), workers in the same firm may optimally face different wage paths, layoff risks, and responses to productivity shocks.

Lastly, it relates to models that use match heterogeneity to generate layoffs, such as \textcite{berger2011}, \textcite{menzio2011}, and \textcite{gregory2021}. In \textcite{menzio2011} and \textcite{gregory2021}, matches dissolve when they become sufficiently unproductive. \textcite{berger2011} use match heterogeneity to explain countercyclical labor productivity and jobless recoveries. The cycle is similar here: firms grow “fat” in high-productivity states and cleanse their workforce when negative shocks arrive. My contribution adapts this mechanism to the wage–firing trade-off: unlike \textcite{berger2011}, firms are free to change workers’ wages but choose not to because, after layoffs, remaining workers are sufficiently valuable to retain.

Empirically, this paper relates to a set of studies measuring wage pass-through and, separately, layoffs. My average pass-through estimates align with \textcite{souchier2022} and with estimates from other countries (\textcite{guvenen2017}, \textcite{guiso2020}). On layoffs, \textcite{buhai2014} show that last-in workers are typically first out; both my empirical results and my model are consistent with this pattern. Beyond tenure and seniority profiles of wages and layoffs, I document novel differences in pass-through by tenure. The cross-firm link between layoffs and pass-through is consistent with recent structural work by \textcite{enrlich2024}.
